\cotefile{156-9}
\scanpage{42}
\mspage{38}
% archive sheet number: 41

y jouant le rôle des $\mathrm{Inf}$ quelconques, contenant $\varnothing$ (plus petit élément) et $\unsure{\mathcal A|A}$ (plus grand élément). On y a une opération $\unsure{\mathrm{cosupp}^{\circ}}$ ou $C$, qui est une anti-involution. Sa donnée \deleted{implique} (p.ex) implique l'existence des $\mathrm{Sup}$ quelconques dans $\Sigma_{\unsure{\mu|M|U}}$ (anti $\unsure{\mathrm{Inf}|\mathrm{Sup}}$) et $C$ échange $\mathrm{Inf}$ et $\mathrm{Sup}$, plus petit et plus grand élément. La relation $\mathrel{|\circ|}$ \deleted{\unsure{se donne par}} \inserted{s'étend} aux éléments de $\Sigma_\mu$, par
\[
S\mathrel{|\circ|}S',\quad \overset{\inserted{\unsure{\mathrm{Not}|\mathrm{def}}}}{\Longleftrightarrow}\quad \forall x\in S,x'\in S',\ \text{on a }x\mathrel{|\circ|}x'.
\]
On a
\[
S\mathrel{|\circ|}S'\Longleftrightarrow S\subset CS'\Longleftrightarrow S'\subset CS.
\]
On a
\[
S\mathrel{|\circ|}S'\Longrightarrow S\cap S'=\varnothing
\]
mais l'inverse n'est pas vrai en général.

% margin note written sideways
% slanted
\margin{Structure d'un \uline{\unsure{quasi-radicié|quasi-radicoïde}}\par
Spatios :\par
a) relation d'ordre\par
b) application $C$\par
\unsure{Alors que} : $1^{\circ})$ $\exists$ plus petit et plus grd élément\par
$2^{\circ})$ $C$ est une anti-involution\par
$3^{\circ})$ \deleted{$\unsure{\forall x\in\Sigma}$} $x\wedge Cx=\varnothing_\Sigma$\par
$(\Longleftrightarrow x\vee Cx=1_\Sigma)$ \unsure{Description|Dans l'option}\par
$\longrightarrow\ \unsure{\mathrm{Figel}\,\mathrm{ord}|\mathrm{Figbord}}(\Sigma)$.}

La donnée d'un magasin \unsure{abstrait|étalon} $\mathcal A$, soit $\unsure{\mathscr F|\mathcal F|\mathbb F},$ (comme ``figures''), \deleted{permet} plus la donnée de $\trianglelefteq$, permet de définir dans $\Sigma_{\mathcal A}$ une partie, formée des \uline{supports constructibles}
\scanpage{43}
\mspage{39}
% archive sheet number: 42

\[
\unsure{\Sigma\mathrm{cons}_{\mathcal A}|\Sigma_{\mathrm{cons}_{\mathcal A}}}
\]
Elle n'est \deleted{\unsure{en}} en général stable ni par $\mathrm{Sup}$ ni par $\mathrm{Inf}$ ni par $C$. Mais il y a un axiome (``axiome des supports'') qui tient lieu d'une telle compatibilité, sur laquelle nous allons avoir à revenir --- et à vérifier \inserted{encore} cet axiome, en faisant intervenir cette fois aussi la structure $\unsure{\overset{\circ}{\ll}|\overset{\circ}{\Subset}}$.

\deleted{Pour l'instant, je suis \unsure{resté}} \inserted{Du point de vue notation} pour $X\in\mathcal A$, je dénote par
\[
-\underset{\substack{\Vert\\\inserted{\mathrm{supp}^{\circ}(X)}}}{X^{\circ}}\in\Sigma\mathrm{cons}_{\mathcal A}
\]
l'élément $\mathrm{supp}^{\circ}X$, qui peut être décrit comme le plus petit $S\in\Sigma_{\mathcal A}$ (donc aussi de $\Sigma\mathrm{cons}_{\mathcal A}$) tel que $X\in S$. C'est une explicitation possible de l'intuition de \unsure{la|l'} \uline{\unsure{Art|Aut} ouvert} associée à un Art donné $X$.
\scanpage{44}
\mspage{40}
% archive sheet number: 43

Je vais maintenant relier la relation $\unsure{\overset{\circ}{\ll}|\overset{\circ}{\Subset}}$ \inserted{(ou $\ll$)}, qui n'est pas encore intervenue dans ces considérations de supports (si ce n'est pour dire que les $S\in\Sigma_{\mathcal A}$ sont des parties $\overset{\circ}{\ll}$-fermées de $\mathcal A$), à ces notions de support. Tout d'abord, pour $\forall X\in\mathcal A$, considérons \deleted{$\unsure{X\leftarrow}$} l'ens
\[
\unsure{\mathrm{Dep}|\mathrm{D\acute{e}p}}(X)=\{\deleted{\unsure{X}}X^{\prime\circ}\mid X'\in\widetilde X\}\subset\unsure{\Sigma\mathrm{cons}|\Sigma_{\mathrm{cons}}}\subset\Sigma
\]
(``\uline{\unsure{déployement}}\deleted{'' de $X$} (``spatial'') de $X$).

On pourrait le définir pour $\unsure{H\subset|\forall}F\in\unsure{\mathscr F|\mathcal F|\mathbb F}(\mathcal A)$, mais pour l'instant on ne s'y intéresse que pour les $X\in\mathcal A$. On a

$\mathrm{Dep}(X)$ admet un plus grand élément, $X^{\circ}$, et est formé d'éléments \inserted{$\ne\varnothing_\Sigma$,} \unsure{mutuellement} disjoints ($\mathrel{|\circ|}$) de $\Sigma\mathrm{cons}$. De plus, l'application $\widetilde X\mathbin{-}\mathrm{Dep}(X)$ est bijective, et $\mathrm{Dep}(X)$ hérite de $\widetilde X$ une structure d'ordre. On considérera tjrs $\mathrm{Dep}(X)$ comme une partie de $\Sigma$, munie de \deleted{une} \inserted{sa} relation d'ordre, \unsure{cela aussi}, comme une figure (élémentaire) ensembliste dans $\Sigma$, ``\uline{séparée}'' i.e. le préordre sur son support est séparé.
\scanpage{45}
\mspage{41}
% archive sheet number: 44

Supposons maintenant qu'on ait
\[
X\ll Y\qquad (X,Y\in\mathcal A)
\]
on a alors une application croissante canonique
\[
\unsure{\varphi|\phi}:\widetilde X\longrightarrow\widetilde Y
\]
telle que, pour \deleted{tout} $X'\in\widetilde X$, on ait
\[
X'\unsure{\overset{\circ}{\ll}|\overset{\circ}{\Subset}}\deleted{\varphi(X')\overset{\mathrm{def}}{=}}Y'\overset{\mathrm{def}}{=}\varphi(X').
\]
En termes de \unsure{déployements}, on peut interpréter $\varphi$ par une application croissante $\unsure{\psi|\tau}$
\[
\psi:\mathrm{D\acute{e}p}(X)\longrightarrow\mathrm{D\acute{e}p}(Y)
\]
satisfaisant
\[
\forall S\in\mathrm{D\acute{e}p}\,X,\qquad S\leq\psi(S)
\]
(ici on a noté $\leq$ la relation $\subset$ dans $\Sigma$).

Cette application $\psi$ est uniquement déterminée par cette propriété, car \deleted{si} \deleted{\unsure{peut avoir|pour avoir}}
\[
S\leq T,T',\qquad\text{avec }T,T'\in\mathrm{Dep}(Y)
\]
on a $T=T'$ (et ceci, pour peu seulement que $S$ soit un élément de $\Sigma$, $\ne\varnothing_\Sigma$). En effet, si $T\ne T'$, on a $T\mathrel{|\circ|}T'$,
\scanpage{46}
\mspage{42}
% archive sheet number: 45

donc $T\cap T'=\varnothing$, donc $S\leq\varnothing$ donc $S=\varnothing$.

Cela impose donc une condition \inserted{nécessaire} en termes des \deleted{supports} \inserted{\unsure{déployements}}, pour qu'on ait
\[
X\ll\deleted{\unsure{X}}Y,
\]
c'est la condition suivante :
\[
\begin{aligned}
&\forall S\in\unsure{\mathrm{Dep}|\mathrm{D\acute{e}p}}\,X,\quad\exists T\in\mathrm{Dep}\,Y,\ \inserted{\text{(néc. unique, \unsure{not|tel} }T=\unsure{\psi|\tau}(S)\text{)}}\ \text{avec}\\
&S\leq T.
\end{aligned}
\tag*{\uline{Raff}}
\]
De plus, l'application
\[
\psi:\mathrm{Dep}(X)\mathbin{-}\mathrm{Dep}(Y)
\]
est croissante.

De plus, on a $X\unsure{\overset{\circ}{\ll}|\overset{\circ}{\Subset}}Y$ \unsure{sss|ssi} l'application $\psi$ précédente transforme le plus grand élément $X^{\circ}$ en le plus grand élément $Y^{\circ}$, i.e. si et seulement si on a
\[
\deleted{\unsure{X^{\circ}\ll}}\qquad X^{\circ}\leq Y^{\circ}
\]
(\uline{NB} \unsure{il est} entendu cependant qu'on sait déjà que $X\ll Y$\dots).

% margin note written sideways
% slanted
\margin{Si l'on a $X\trianglelefteq Y$, alors $\psi:\mathrm{Dep}\,X\longrightarrow\mathrm{Dep}\,Y$ est une inclusion d'ensembles ordonnés $\mathrm{Dep}(X)\subset\mathrm{Dep}(Y)$, et $\mathrm{Dep}(X)$ est une partie \unsure{ordonnée} \uline{fermée} de $\mathrm{Dep}\,Y$.}

\uline{Définition}. On dit qu'un \deleted{\unsure{magasin}} \inserted{atelier} est \uline{spatial}, si la condition nécessaire précédente Raff, sur $\mathrm{D\acute{e}p}\,X$ et $\mathrm{D\acute{e}p}\,Y$, pour que l'on ait $X\ll Y$,
\scanpage{47}
\mspage{43}
% archive sheet number: 46

est aussi suffisante.

Le terme ``spatial'' \deleted{\unsure{signi}} est destiné ici à suggérer que la description de la relation \unsure{primitive} $\unsure{\overset{\circ}{\ll}|\overset{\circ}{\Subset}}$ peut se faire en termes de $\trianglelefteq$, $\mathrel{|\circ|}$, par référence à des notions de nature purement ``spatiales'' (la notion de support, qui \deleted{fait le lien avec la} correspond à la notion ``d'espace occupé'' par une figure\unsure{.\,--|\dots}).

Cela nous met devant la question : étant donné une structure
\[
\deleted{\unsure{\mathcal A}}(\mathcal A,\trianglelefteq,\mathrel{|\circ|}),
\]
\unsure{sous} quelles conditions provient-elle d'une structure d'atelier spatial
\[
(\mathcal A,\trianglelefteq,\overset{\circ}{\ll},\mathrel{|\circ|})\quad ?
\]

La structure de départ $(\mathcal A,\trianglelefteq,\mathrel{|\circ|})^{*}$ permet de construire l'ens ordonné $\Sigma$, avec la structure supplémentaire $C$ (permettant de définir $\mathrel{|\circ|}$ sur $\Sigma$), et l'application :
% margin note written sideways
% slanted
\margin{$*+$ l'axiome \unsure{At 4} qu'il faut \unsure{rappeler} ici, comme \unsure{condition} (i)}
\[
X\longrightarrow\unsure{\mathrm{Dep}|\mathrm{D\acute{e}p}}(X)\qquad \mathcal A\longrightarrow\mathrm{Fig\acute{e}l}^{\inserted{\mathrm{ord}}}(\Sigma)
\]
\deleted{vers} \inserted{où $\unsure{\mathrm{Figel}\,\mathrm{ord}|\mathrm{Fig\acute{e}l}^{\mathrm{ord}}}(\Sigma)$ désigne l'ens.} les figures élémentaires \inserted{\unsure{ens. ord|ens. sép}} de $\Sigma$ qui
\scanpage{48}
\mspage{44}
% archive sheet number: 47

sont compatibles avec la relation \deleted{\unsure{d'ordre}} \inserted{\unsure{(}d'ensemble \deleted{\unsure{ordonné}} support}; \deleted{savoir} i.e. telles que \deleted{$\unsure{\forall x\in A\in}$} $S\in\unsure{\Phi|\varphi}$
\[
\begin{aligned}
\text{a)}\quad&S\in\Phi\Longrightarrow S\ne\varnothing_\Sigma\\
\text{b)}\quad&S,T\in\Phi,\ S\ne T\Longrightarrow S\leq CT\quad\text{i.e. }T\leq CS.
\end{aligned}
\]
\deleted{La première condition nécessaire, qui exprime la}

La condition
\[
\mathrm{D\acute{e}p}(X)\unsure{\overset{\circ}{\ll}|\overset{\circ}{\Subset}}\mathrm{D\acute{e}p}(Y)
\]
est une relation de préordre sur $\mathcal A$, que je noterai
\[
X\overset{\circ}{\ll}Y
\]
La première condition, c'est que cette relation soit une \uline{relation d'ordre}, i.e. \deleted{que}

(ii) Soient $X,Y\in\mathcal A$, tels que $\unsure{\mathrm{Dep}|\mathrm{D\acute{e}p}}\,X=\mathrm{Dep}\,Y$, alors $X=Y$.

Ceci fait, on trouve une structure $(\mathcal A,\trianglelefteq,\overset{\circ}{\ll},\mathrel{|\circ|})$, satisfaisant déjà At 1 a) b) c). Je voudrais que la relation $\ll$ qu'elle définit soit celle \deleted{de} \unsure{déduite|décrite} par $\ll$ dans $\unsure{\mathrm{Figel}\,\mathrm{ord}|\mathrm{Fig\acute{e}l}\,\mathrm{ord}|\mathrm{Fig\acute{e}l}^{\mathrm{ord}}}(\Sigma)$.

\deleted{Cela signifie que si $\mathrm{Dep}(X)$} Supposons donc \deleted{\unsure{De}}
\[
\mathrm{D\acute{e}p}(X)\ll\mathrm{Dep}(Y)
\]
donc
\scanpage{49}
\mspage{45}
% archive sheet number: 48

\[
\unsure{\mathrm{Dep}|\mathrm{D\acute{e}p}}\,X\unsure{\overset{\circ}{\ll}|\overset{\circ}{\Subset}}
\underset{\begin{gathered}
\rule{0.3pt}{1em}\\
\inserted{\begin{gathered}
\text{sous-figure}\\
\text{élémentaire de }Y,\text{ définie}\\
\text{par }\unsure{\psi|\tau}(X^{\circ})
\end{gathered}}
\end{gathered}}{\unsure{\Phi|\varphi}}
\trianglelefteq\!\leftarrow\mathrm{Dep}(Y)
\]

on sait que $\Phi$ sera de la forme
\[
\deleted{\unsure{\mathrm{De}}}\qquad\Phi=\mathrm{D\acute{e}p}(Y'),\qquad Y'\trianglelefteq Y,
\]
et on aura \deleted{alors} bien alors, par définition $X\overset{\circ}{\ll}Y'$, \unsure{oK|OK|ok}.

\uline{At 2} a) immédiat par \unsure{a)|(ii)}

b) Soit
\diagram{\begin{tikzcd}[ampersand replacement=\&,column sep=large,row sep=large]
X\arrow[r,phantom,"{\overset{\circ}{\ll}}"{description}]\arrow[d,phantom,"{\trianglelefteq}"{description,rotate=90}]\&Y\\
X'\&
\end{tikzcd}}
, prouvons $\exists Y'\in\widetilde Y$ avec $X'\overset{\circ}{\ll}Y'$. Immédiat.

\uline{At 3} On voit que la relation $X\mathrel{|\circ|}Y$ se \deleted{\unsure{transpose par des}} \inserted{traduit sur} $\mathrm{Dep}\,X$, $\mathrm{Dep}(Y)$ par

\deleted{plus} Si $S$ est le plus grand élément de $\mathrm{Dep}\,X$\par
$-T\quad\rule{7cm}{0.3pt}\quad\mathrm{Dep}\,Y$,\par
on a $S\mathrel{|\circ|}T$ dans $\Sigma$.

Et l'axiome At 3 résulte de la relation évidente dans $\Sigma$
\[
S\mathrel{|\circ|}T,\quad S'\leq S,T'\leq T\Longrightarrow S'\mathrel{|\circ|}T'
\]

\uline{At 4} \deleted{Evident.} Ne fait pas intervenir $\overset{\circ}{\ll}$, \deleted{déjà connu.} il faut bien sûr l'imposer d'avance.
